\NativeTitle{53}{第六单元素养评价（B）}{（小数的意义与性质）}{90}{5}{一 & 二 & 三 & 四 & 五}{ &  &  &  & }
\NativeSlot{282}{36}{%
\ExamGroup{一、填空。{\normalfont\kaishu\fontsize{10.8bp}{14bp}\selectfont （每空1分，共23分）}}

% p053_s01

\ExamQuestion{1.}{5}%
和1相邻的一位小数分别是\mbox{（\hspace{1.8em}）}和\mbox{（\hspace{1.8em}）}。\par
}
\NativeSlot{332.0}{43}{%
% p053_s02

\ExamQuestion{2.}{5}%
科学家们发现，当外界环境温度达到人体正常温度与0.618的乘积时，人会感到最舒服。0.618的计数单位是\mbox{（\hspace{4.5em}）}，它含有\mbox{（\hspace{1.8em}）}个这样的计数单位，如果改写成以万分之一为计数单位的数，这个数是\mbox{（\hspace{1.8em}）}。\par
}
\NativeSlot{409}{67}{%
% p053_s03

\ExamQuestion{3.}{5}%
图中两个“4”表示的意义不同，这两个“4”相差\mbox{（\hspace{1.8em}）}。\par
\BookFigure{figures/p053_s03-01.tex}
}
\NativeSlot{521.0}{36}{%
% p053_s04

\ExamQuestion{4.}{5}%
在0.03、0.003、0.030这三个数中，计数单位相同的是\mbox{（\hspace{1.8em}）}和\mbox{（\hspace{1.8em}）}；大小相等的是\mbox{（\hspace{1.8em}）}和\mbox{（\hspace{1.8em}）}。\par
}
\NativeSlot{571}{36}{%
% p053_s05

\ExamQuestion{5.}{5}%
一个数由3个十，8个十分之一和5个千分之一组成，这个数写作\mbox{（\hspace{4.4em}）}，读这个数时，有\mbox{（\hspace{1.8em}）}个零不必读出来。\par
}
\NativeSlot{622.0}{36}{%
% p053_s06

\ExamQuestion{6.}{5}%
甲数是3.8，乙数是38，在它们的末尾都填上两个零，这时乙数是甲数的\mbox{（\hspace{1.8em}）}倍。\par
}
\NativeSlot{673.5}{36}{%
% p053_s07

\ExamQuestion{7.}{5}%
一个三位小数，各个数位上的数字之和是9，且各个数位上的数字都不相同，符合条件的三位小数最小是\mbox{（\hspace{1.8em}）}。\par
}
\NativeSlot{723}{36}{%
% p053_s08

\ExamQuestion{8.}{5}%
在括号里填上合适的单位，使等式成立。45\mbox{（\hspace{1.8em}）}＝0.45\mbox{（\hspace{1.8em}）}；3540\mbox{（\hspace{1.8em}）}＝3.540\mbox{（\hspace{1.8em}）}。\par
}
\NativeSlot{774.5}{36}{%
% p053_s09

\ExamQuestion{9.}{5}%
10千克花生可以榨4.5千克花生油，1吨花生可以榨\mbox{（\hspace{1.8em}）}千克花生油，合\mbox{（\hspace{1.8em}）}吨花生油。（第二空接下页）\par
}
